\documentclass[11pt,a4paper]{article}

% ---------- Packages ----------
\usepackage[margin=1in]{geometry}
\usepackage[dvipsnames]{xcolor}
\usepackage{amsmath, amssymb, amsthm, mathtools}
\usepackage{algorithm}
\usepackage{algpseudocode}
\usepackage[colorlinks=true,
            linkcolor=blue!50!black,
            urlcolor=blue!50!black,
            citecolor=blue!50!black]{hyperref}

% ---------- Theorem environments ----------
\theoremstyle{plain}
\newtheorem{theorem}{Theorem}[section]
\newtheorem{proposition}[theorem]{Proposition}
\newtheorem{lemma}[theorem]{Lemma}
\newtheorem{corollary}[theorem]{Corollary}

\theoremstyle{definition}
\newtheorem{definition}[theorem]{Definition}

\theoremstyle{remark}
\newtheorem{remark}[theorem]{Remark}
\newtheorem{example}[theorem]{Example}

% ---------- Math shortcuts ----------
\newcommand{\R}{\mathbb{R}}
\newcommand{\N}{\mathbb{N}}
\newcommand{\E}{\mathbb{E}}
\newcommand{\Var}{\operatorname{Var}}
\newcommand{\Cov}{\operatorname{Cov}}
\newcommand{\Corr}{\operatorname{Corr}}
\newcommand{\F}{\mathcal{F}}
\newcommand{\Prob}{\mathbb{P}}
\newcommand{\Q}{\mathbb{Q}}
\newcommand{\Tcal}{\mathcal{T}}
\newcommand{\dd}{\mathop{}\!\mathrm{d}}
\newcommand{\eqdef}{\overset{\mathrm{def}}{=}}
\newcommand{\indic}{\mathbf{1}}
\newcommand{\esssup}{\operatorname*{ess\,sup}}

% ---------- Document metadata ----------
\title{American Options by Longstaff-Schwartz:\\
       Regression-Based Optimal Stopping with a Binomial-Tree Validator}
\author{Quant Finance Portfolio --- Phase 2, Block 6}
\date{\today}

\begin{document}
\maketitle

\begin{abstract}
The American option pricing problem is fundamentally different from
its European counterpart: the holder can exercise at any time up to
maturity, so the price is the value function of a continuous-time
optimal stopping problem. No closed form is available except in
degenerate cases (zero rate, infinite maturity). This block presents
the Longstaff-Schwartz (2001) regression-based Monte Carlo algorithm,
which discretises the exercise dates to a Bermudan grid and replaces
the conditional expectation in the dynamic-programming Bellman
equation by a cross-sectional regression of continuation values onto
a finite basis of polynomials of the underlying. We use Laguerre
polynomials of degree up to three for the regression basis. The
algorithm is validated against an independent Cox-Ross-Rubinstein
binomial tree implementation, also from scratch, which serves as the
ground truth in the limit of small step size.
\end{abstract}

\tableofcontents

\section{Optimal stopping and the American option}
\label{sec:setting}

\subsection{The American put}

An \emph{American put} on a single non-dividend-paying asset $S$ with
strike $K$ and maturity $T$ pays its holder $\max(K - S_\tau, 0)$ at
any chosen time $\tau \in [0, T]$, at the holder's discretion. Under
the risk-neutral measure $\Q$ with constant rate $r$, the time-zero
price is
\begin{equation}
V_0 \;=\; \sup_{\tau \in \Tcal_{[0, T]}}
            \E^{\Q}\!\left[ e^{-r\tau}\, (K - S_\tau)^+ \right],
\label{eq:american-price}
\end{equation}
where $\Tcal_{[0, T]}$ is the set of stopping times of the natural
filtration $\F = \{\F_t\}$ taking values in $[0, T]$, and
$x^+ \eqdef \max(x, 0)$.

The supremum is attained at an \emph{optimal stopping time} $\tau^\star$.
The map $t \mapsto S_t^\star$ separating the continuation region
$\{S > S_t^\star\}$ from the exercise region $\{S \le S_t^\star\}$ is
the \emph{optimal exercise boundary}; it is a non-trivial function of
$t$ for the American put and admits no closed form.

\subsection{Why no Black-Scholes formula}

For the European put, the holder is committed to exercising at $T$,
so the option price is just the discounted expectation
$e^{-rT}\,\E^\Q[(K - S_T)^+]$, which integrates against the lognormal
density of $S_T$ to produce the Black-Scholes formula. The American
case introduces an entirely different kind of object: the supremum
over $\tau$ in~\eqref{eq:american-price} is a sup over an
infinite-dimensional set of measurable strategies. There is no
elementary integral that captures it.

The classical machinery for handling \eqref{eq:american-price} is the
\emph{Snell envelope} construction.

\subsection{The Snell envelope}

For a discounted reward process $X_t = e^{-rt} (K - S_t)^+$ adapted
to $\F$, the Snell envelope $U$ is defined as the smallest
supermartingale dominating $X$:
\begin{equation}
U_t \;=\; \esssup_{\tau \in \Tcal_{[t, T]}} \E^\Q[X_\tau \mid \F_t].
\end{equation}
Under standard regularity, $U_t$ exists and is the unique solution to
the recursive equation
\begin{equation}
U_t \;=\; \max\!\bigl(\, X_t,\, \E^\Q[U_{t+\dd t} \mid \F_t]\, \bigr),
\qquad U_T = X_T.
\label{eq:snell}
\end{equation}
The price of the American option is $V_0 = U_0$, and the optimal
stopping time is $\tau^\star = \inf\{ t \ge 0 : U_t = X_t \}$ --- the
first time the envelope touches the immediate-exercise reward.

Equation~\eqref{eq:snell} is a backward dynamic programming
recursion. In continuous time it leads to a free-boundary PDE (the
Black-Scholes PDE with an obstacle); in discrete time, which we
adopt below, it gives a clean numerical scheme.

\section{Bermudan approximation and the Bellman equation}
\label{sec:bermudan}

\subsection{Discretising the exercise dates}

A \emph{Bermudan} option is the discrete analogue: exercisable only at
a finite set of pre-specified dates $0 = t_0 < t_1 < \cdots < t_N = T$.
We take an equispaced grid $t_k = kT/N$ for some large $N$. As $N \to
\infty$, the Bermudan price converges to the American price
monotonically from below~\cite[Section 8.3]{Glasserman}.

Let $S_k \eqdef S_{t_k}$. The Bermudan put value at exercise date
$t_k$, given $S_k = s$, satisfies the Bellman recursion:
\begin{equation}
\boxed{
\begin{aligned}
V_N(s) \;&=\; (K - s)^+, \\[0.4em]
V_k(s) \;&=\; \max\!\bigl( \underbrace{(K - s)^+}_{\text{intrinsic}},\;
\underbrace{e^{-r\Delta t}\,\E^\Q[V_{k+1}(S_{k+1}) \mid S_k = s]}_{\text{continuation value }C_k(s)} \bigr),
\end{aligned}
}
\label{eq:bellman}
\end{equation}
for $k = N - 1, N - 2, \ldots, 0$, with $\Delta t = T/N$. Then
$V_0(S_0)$ is the Bermudan put price.

The Markov property of GBM under $\Q$ guarantees that the conditional
expectation $\E^\Q[V_{k+1}(S_{k+1}) \mid \F_k]$ depends on $\F_k$
only through $S_k$, so the continuation value $C_k(s)$ is a function
of one scalar variable. This is essential for the regression-based
approach below.

\subsection{The optimal exercise rule}

Equation~\eqref{eq:bellman} also gives the optimal exercise rule:
exercise at the first $k$ such that
\[
(K - S_k)^+ \;\ge\; C_k(S_k).
\]
At each path, the rule depends only on the comparison between the
known intrinsic value and the unknown continuation value. The whole
challenge is computing $C_k(s)$.

\subsection{Why pure Monte Carlo fails directly}

A naive attempt to estimate $C_k(s)$ would require, at each path
$i$ and each time $k$, simulating a sub-tree of paths from
$(t_k, S_k^{(i)})$ to $(t_{k+1}, S_{k+1})$ and averaging the
continuation values $V_{k+1}(S_{k+1})$. This is a nested simulation
of cost $O(n^2 N)$ for $n$ paths and $N$ time steps --- prohibitive
for any realistic $N$ and $n$. Worse, $V_{k+1}$ is itself defined by
nested simulation, so the cost compounds geometrically with $k$.

Longstaff and Schwartz~\cite{LongstaffSchwartz2001} solved this in a
single line of insight: \emph{the conditional expectation is a
function of $S_k$, so we can fit it by cross-sectional regression
across the existing paths.}

\section{The Longstaff-Schwartz algorithm}
\label{sec:lsm}

\subsection{The regression idea}

At time $t_k$, the unknown continuation value $C_k(\cdot): \R_+ \to
\R$ is approximated by a linear combination of $M$ basis functions
$\psi_1, \psi_2, \ldots, \psi_M$ of the underlying:
\begin{equation}
C_k(s) \;\approx\; \widehat C_k(s) \;\eqdef\;
\sum_{m=1}^M \beta_m^{(k)}\, \psi_m(s).
\end{equation}
The coefficients $\beta^{(k)} = (\beta_1^{(k)}, \ldots, \beta_M^{(k)})$
are chosen by ordinary least squares against the simulated cash flows
of the optimal strategy from $t_{k+1}$ onwards.

Concretely: simulate $n$ independent paths
$(S_0, S_1^{(i)}, \ldots, S_N^{(i)})_{i = 1, \ldots, n}$ of the GBM
under $\Q$. Process the time grid \emph{backwards}. At step $k$, we
have already determined, for each path $i$, the discounted cash flow
$Y^{(i)}$ that the optimal strategy from time $t_{k+1}$ onwards
produces, evaluated at time $t_k$ (i.e., discounted back to $t_k$).
Regress these $\{Y^{(i)}\}_i$ on the $M$-vector of basis values
$\{\psi_m(S_k^{(i)})\}_m$:
\begin{equation}
\beta^{(k)} \;=\; (\Psi_k^\top \Psi_k)^{-1} \Psi_k^\top Y,
\label{eq:ols}
\end{equation}
where $\Psi_k$ is the $n \times M$ design matrix with entries
$(\Psi_k)_{i, m} = \psi_m(S_k^{(i)})$. The regression produces an
estimated continuation value $\widehat C_k(S_k^{(i)})$ at each path.

Then update the cash flow array: at any path $i$ where the intrinsic
value $(K - S_k^{(i)})^+$ exceeds $\widehat C_k(S_k^{(i)})$, the
optimal action is to exercise now, so $Y^{(i)}$ is replaced by
$(K - S_k^{(i)})^+$. Otherwise $Y^{(i)}$ is discounted back one step:
$Y^{(i)} \to e^{-r\Delta t} Y^{(i)}$.

After the recursion reaches $k = 0$, the price estimator is the
sample mean of the $\{Y^{(i)}\}_i$ discounted to $t_0$:
\begin{equation}
\widehat V_0 \;=\; \frac{1}{n} \sum_{i=1}^n Y^{(i)}.
\label{eq:price-estimator}
\end{equation}

\subsection{The in-the-money restriction}

A critical detail of the original LS paper: the regression
\eqref{eq:ols} should be performed only on the subset of paths that
are \emph{in the money} at $t_k$, i.e., for which $K - S_k^{(i)} > 0$.
Out-of-the-money paths contribute zero intrinsic value, so the
exercise vs.\ continuation decision is moot at those paths (always
continue). Including them in the regression dilutes the signal: it
forces the polynomial fit to behave reasonably where it doesn't
matter, at the cost of accuracy where it does matter.

This is one of those small algorithmic choices that makes a large
practical difference. We follow it in the implementation.

\subsection{Algorithm pseudocode}

\begin{algorithm}[H]
\caption{Longstaff-Schwartz for the American put}
\label{alg:lsm}
\begin{algorithmic}[1]
\Require Spot $S_0$, params $(K, r, \sigma, T)$, exercise grid $N$, paths $n$, basis $\{\psi_m\}_{m=1}^M$.
\State Simulate matrix $\mathbf{S} \in \R^{n \times (N+1)}$ of GBM paths.
\State Initialise cash flow $Y^{(i)} \gets (K - S_N^{(i)})^+$ for all $i$.
\For{$k = N - 1$ \textbf{downto} $1$}
  \State Discount cash flows: $Y^{(i)} \gets e^{-r\Delta t} Y^{(i)}$.
  \State Find ITM paths: $\mathcal{I}_k = \{i : K - S_k^{(i)} > 0\}$.
  \State Build design $\Psi_k \in \R^{|\mathcal{I}_k| \times M}$, target $Y_{\mathcal{I}_k}$.
  \State Solve $\beta^{(k)} = (\Psi_k^\top \Psi_k)^{-1} \Psi_k^\top Y_{\mathcal{I}_k}$.
  \State Compute $\widehat C_k^{(i)} = \sum_m \beta_m^{(k)} \psi_m(S_k^{(i)})$ for $i \in \mathcal{I}_k$.
  \State Determine exercise set $\mathcal{E}_k = \{i \in \mathcal{I}_k : (K - S_k^{(i)}) \ge \widehat C_k^{(i)}\}$.
  \State For $i \in \mathcal{E}_k$: replace $Y^{(i)} \gets K - S_k^{(i)}$.
\EndFor
\State Final discount: $Y^{(i)} \gets e^{-r\Delta t} Y^{(i)}$ for all $i$.
\State \Return $\widehat V_0 = \frac{1}{n}\sum_i Y^{(i)}$ and standard MC half-width.
\end{algorithmic}
\end{algorithm}

\paragraph{Note on time $k = 0$.} The algorithm exits the loop at
$k = 1$ and applies one final discount. This is because $S_0$ is
deterministic, so the exercise-or-continue decision at $t_0$ is just a
comparison of the immediate intrinsic value $(K - S_0)^+$ against the
sample mean $\widehat V_0$ produced by the algorithm. We return the
maximum of the two. (For the canonical at-the-money put this is
moot, since $(K - S_0)^+ = 0$.)

\section{The Laguerre basis}
\label{sec:laguerre}

\subsection{Definition}

The Laguerre polynomials $\{L_n\}_{n \ge 0}$ are the orthogonal
polynomials on $[0, \infty)$ with respect to the weight
$w(x) = e^{-x}$. They are defined by the recurrence
\begin{equation}
L_0(x) = 1, \quad L_1(x) = 1 - x,
\quad
(n + 1) L_{n+1}(x) = (2n + 1 - x) L_n(x) - n\, L_{n-1}(x).
\end{equation}
Closed forms for the first few:
\begin{align}
L_0(x) &= 1, \\
L_1(x) &= 1 - x, \\
L_2(x) &= \tfrac{1}{2}(x^2 - 4x + 2), \\
L_3(x) &= \tfrac{1}{6}(-x^3 + 9x^2 - 18x + 6).
\end{align}
Their orthogonality is
\(
\int_0^\infty L_n(x) L_m(x)\, e^{-x}\,\dd x \;=\; \delta_{nm}.
\)

\subsection{Use in LSM}

In our LSM implementation we apply the basis to the normalised
underlying $x = S/K$, with $M = 4$ basis functions:
\begin{equation}
\psi_m(S) = L_{m-1}(S/K), \qquad m = 1, 2, 3, 4.
\end{equation}
The normalisation $S/K$ keeps the argument in a controlled range
(typically in $[0, 1]$ for ITM put paths, since ITM means $S < K$).
On $[0, 1]$, the unweighted Laguerre polynomials are well-behaved and
linearly independent; the design matrix $\Psi_k^\top \Psi_k$ is
positive-definite and Cholesky-factorisable.

\paragraph{Why Laguerre and not monomials?} The original
Longstaff-Schwartz paper uses weighted Laguerre $\tilde\psi_m(x) =
e^{-x/2} L_{m-1}(x)$, which is orthonormal under $\dd x$. The weight
makes the basis bounded uniformly on $[0, \infty)$. Our use of the
unweighted form is acceptable because we restrict regression to ITM
paths (so $x \in [0, 1]$ effectively), where the weight is close to
constant. The numerical behaviour is essentially identical; we choose
unweighted for code simplicity. Plain monomials $\{1, x, x^2, x^3\}$
also work and are commonly used by practitioners; the choice of basis
matters less than the principle of \emph{using one}.

\section{The Cox-Ross-Rubinstein binomial tree validator}
\label{sec:binomial}

The LSM estimator is statistically biased and depends on basis
choice, sample size, and number of exercise dates. To validate it we
need an independent reference price computed by a different method.
The standard choice is the Cox-Ross-Rubinstein (CRR)
\emph{binomial tree}~\cite{CRR1979}, which converges to the true
American price as the tree depth increases.

\subsection{Construction}

Discretise time into $N$ steps of size $\Delta t = T/N$. At each step
the asset goes up by factor $u$ or down by factor $d$, with
risk-neutral probabilities $(p, 1 - p)$:
\begin{equation}
u = e^{\sigma \sqrt{\Delta t}},
\quad
d = 1/u = e^{-\sigma \sqrt{\Delta t}},
\quad
p = \frac{e^{r \Delta t} - d}{u - d}.
\end{equation}
The choice $u \cdot d = 1$ keeps the tree recombining: after $j$ ups
and $k - j$ downs, the asset price is $S_0 \cdot u^{2j - k}$. The
total number of nodes after $N$ steps is $(N + 1)(N + 2)/2$, so the
tree fits in memory for $N$ up to several thousand.

\subsection{Backward induction}

At maturity ($k = N$), the option value at each terminal node is the
intrinsic put value. At earlier nodes, we apply the Bellman equation
exactly (no regression needed):
\begin{equation}
V_k(j) \;=\; \max\!\bigl( K - S_k(j),\; e^{-r \Delta t}
[\,p\, V_{k+1}(j + 1) + (1 - p)\, V_{k+1}(j)\,] \bigr),
\end{equation}
where $V_k(j)$ is the option value at node $j$ of step $k$ and
$S_k(j) = S_0 u^{2j - k}$. The price is $V_0(0)$.

\subsection{Convergence}

The CRR price converges to the true American price as $N \to \infty$
at rate roughly $1/\sqrt N$ (see~\cite[Section 4.5]{Hull}); in
practice $N = 1000$ to $5000$ is sufficient for four-decimal accuracy
on a one-year ATM American put with $\sigma = 0.20$.

\section{Bias and convergence of the LSM estimator}
\label{sec:bias}

\subsection{Two sources of error}

The LSM estimator $\widehat V_0^{\mathrm{LSM}}$ has two distinct
sources of error compared to the true Bermudan price:

\begin{enumerate}
\item \textbf{Stochastic error}: the Monte Carlo sampling
      uncertainty, of order $1/\sqrt n$ in the number of paths. This
      is the standard MC error and is captured by the asymptotic
      half-width.

\item \textbf{Approximation bias}: the regression-based continuation
      value $\widehat C_k$ is an approximation of the true $C_k$.
      Two effects compete:

      \begin{itemize}
      \item \emph{Sub-optimal stopping}: any approximation of
            $C_k$ leads to a stopping rule different from $\tau^\star$,
            which gives a strictly lower expected payoff (because
            $\tau^\star$ is by definition the supremum). This is a
            \textbf{negative bias} (LSM tends to underprice).
      \item \emph{In-sample bias}: when we use the same paths to
            estimate $\widehat C_k$ and to apply the resulting rule,
            the rule is data-mined to look better than it really is
            on those paths. This is a \textbf{positive bias} (LSM
            tends to overprice).
      \end{itemize}
\end{enumerate}

In practice the two biases approximately cancel: the LSM estimator is
near-unbiased for sufficiently many paths and a sufficiently rich
basis. The literature documents agreement with reference values to
better than $1\%$ on the canonical American put problem with $n =
10^5$ and $M = 3$ to $4$ Laguerre basis functions.

\subsection{The unbiased upper bound}

A rigorously unbiased \emph{lower bound} estimator of the price can
be obtained by splitting the simulation into two independent batches:
the first batch is used to fit the regression coefficients, the
second to apply the resulting exercise rule and average the cash
flows. This eliminates the in-sample bias at the cost of a constant
factor of paths.

A separate \emph{upper bound} can be constructed via the
Andersen-Broadie (2004) duality approach, which exploits the
martingale formulation of the optimal stopping problem. We do not
implement either refinement in this block; we rely on the
near-cancellation of the two biases and validate against the
binomial tree.

\section{Empirical validation strategy}
\label{sec:validation}

The validation script implements four tests.

\paragraph{Test 1: binomial tree convergence.} Compute the binomial
American put at $N = 100, 500, 1000, 5000, 10000$. The sequence
should converge to a stable value at four decimal places by
$N = 5000$.

\paragraph{Test 2: LSM matches binomial.} At canonical parameters
($S = K = 100$, $r = 0.05$, $\sigma = 0.20$, $T = 1$), with $N_{LSM} =
50$ exercise dates and $n = 10^5$ paths, the LSM estimator should
agree with the binomial reference (at $N_{tree} = 5000$) within
three half-widths.

\paragraph{Test 3: early exercise premium.} Compute both the European
put (Black-Scholes) and the American put (LSM, plus binomial as
sanity). The American minus European spread is the early-exercise
premium; it must be strictly positive, and large enough to be
distinguishable from MC noise (the canonical ATM case gives a
premium of order $0.1$ on a put of order $5.6$).

\paragraph{Test 4: input validation.} Mechanical.

\section{Discussion: when LSM works and when it does not}
\label{sec:discussion}

LSM is the workhorse method for high-dimensional American options
where binomial and PDE methods are infeasible. Its strength is that
the cost scales linearly in the dimension of the underlying state, in
contrast to the exponential blowup of grid-based methods. For our
single-asset American put it is overkill --- the binomial tree is
faster, more accurate, and entirely deterministic --- but the value
of the implementation is in the methodology: LSM is the standard tool
for path-dependent and multi-asset Americans (basket Americans,
swing options, real options on multi-factor commodities).

\paragraph{When LSM fails or struggles.} Three regimes:

\begin{enumerate}
\item \emph{Deeply out-of-the-money options}: the regression has
      few in-the-money paths, so the continuation function is
      poorly identified. Workaround: use importance sampling to
      bias paths into the money.
\item \emph{Highly path-dependent payoffs}: if the option's value
      depends on the entire history rather than just $S_k$, a
      one-dimensional regression is insufficient. Workaround:
      include path-summary statistics in the basis (the running
      maximum, average, etc.).
\item \emph{Discontinuous payoffs}: digital Americans have
      discontinuous intrinsic values, which destabilise the
      regression near the discontinuity. Workaround: smoothing
      via known regularisations.
\end{enumerate}

For the smooth, single-state American put we treat in this block,
none of these difficulties arise.

\section{Implementation outline}
\label{sec:impl}

The block ships two pricers in a new module \texttt{quantlib.american}
(and the C++ mirror):

\begin{itemize}
\item \texttt{binomial\_american\_put}: CRR validator. No randomness,
      just the lattice recursion. Used as the ground truth.
\item \texttt{lsm\_american\_put}: Longstaff-Schwartz with Laguerre
      basis $M = 4$, in-the-money restriction, exact GBM stepping.
\end{itemize}

The C++ side adds a file-local helper \texttt{cholesky\_solve} that
factorises the small $M \times M$ normal-equations matrix and
back-substitutes for $\beta^{(k)}$. For our $M = 4$ basis, this is a
direct rather than iterative solve; the matrix is small,
positive-definite, and well-conditioned in practice.

\section{Bibliographic notes}
\label{sec:biblio}

The Longstaff-Schwartz algorithm is from~\cite{LongstaffSchwartz2001}.
Glasserman~\cite[Chapter 8]{Glasserman} gives a careful textbook
treatment with full convergence theory. The Cox-Ross-Rubinstein
binomial tree is from~\cite{CRR1979}; modern textbook treatment in
Hull~\cite[Chapter 21]{Hull}. Andersen and Broadie's upper-bound
construction is in~\cite{AndersenBroadie2004}.

The treatment of optimal stopping via Snell envelopes is classical;
see e.g.\ Karatzas and Shreve~\cite{KaratzasShreve} for a rigorous
measure-theoretic exposition.

\begin{thebibliography}{99}

\bibitem{LongstaffSchwartz2001}
F.\ A.\ Longstaff and E.\ S.\ Schwartz,
\emph{Valuing American options by simulation: a simple least-squares approach},
Review of Financial Studies, 14(1):113--147, 2001.

\bibitem{CRR1979}
J.\ C.\ Cox, S.\ A.\ Ross, and M.\ Rubinstein,
\emph{Option pricing: a simplified approach},
Journal of Financial Economics, 7(3):229--263, 1979.

\bibitem{AndersenBroadie2004}
L.\ Andersen and M.\ Broadie,
\emph{A primal-dual simulation algorithm for pricing multidimensional American options},
Management Science, 50(9):1222--1234, 2004.

\bibitem{Glasserman}
P.\ Glasserman,
\emph{Monte Carlo Methods in Financial Engineering},
Springer, 2004.

\bibitem{Hull}
J.\ C.\ Hull,
\emph{Options, Futures, and Other Derivatives},
10th edition, Pearson, 2018.

\bibitem{KaratzasShreve}
I.\ Karatzas and S.\ E.\ Shreve,
\emph{Methods of Mathematical Finance},
Springer, 1998.

\end{thebibliography}

\end{document}
